\documentclass[aspectratio=169,10pt]{beamer}
\usepackage{srlecture}
% Shared settings for the three QIPPG103 decks. Change the game URL here once.
\newcommand{\gameurl}{https://ebinrajiiit.github.io/ml_games/}
\newcommand{\gameurlshort}{ebinrajiiit.github.io/ml\_games}
\author{Dr.\ Sara Renjit}
\institute{Indian Institute of Information Technology Kottayam}
\newcommand{\coursename}{QIPPG103 Machine Learning: Principles and Practices}

\title{LSTM, GRU and Bidirectional RNNs}
\renewcommand{\lecturelabel}{QIPPG103 \textperiodcentered\ Lecture 21 \textperiodcentered\ LSTM, GRU, BiLSTM}

\begin{document}

\begin{frame}[plain]
  \begin{tikzpicture}[remember picture,overlay]
    \node[anchor=north east] at ([xshift=-1.2em,yshift=-1.0em]current page.north east) {\includegraphics[height=1.25cm]{iiitk.png}};
  \end{tikzpicture}
  \vskip3.2em
  {\small\scshape\color{inksoft}\coursename\par}
  \vskip1.2em
  {\color{navy}\fontsize{28}{32}\selectfont LSTM, GRU and Bidirectional RNNs\par}
  \vskip0.5em
  {\large\itshape Giving the network a choice about what to remember\par}
  \vskip0.6em{\color{rulegrey}\rule{0.45\textwidth}{0.4pt}}\par\vskip0.8em
  {\normalsize Dr.\ Sara Renjit\par}
  {\small\color{inksoft}Indian Institute of Information Technology Kottayam\par}
  \vskip0.4em
  {\small\color{inksoft}Lecture 21 of 27 \textperiodcentered\ offline session \textperiodcentered\ 90 minutes\par}
  \begin{tikzpicture}[remember picture,overlay]
    \node[anchor=south east] (fl) at ([xshift=-1.4em,yshift=1.2em]current page.south east) {\includegraphics[height=2.1cm]{factsh_logo.png}};
    \node[anchor=east,align=right,font=\scriptsize,text=inksoft] at ([xshift=-0.3em]fl.west) {a session from\\{\color{petrol}\factshfont\bfseries\normalsize FACTS-H Lab}};
  \end{tikzpicture}
\end{frame}

\begin{frame}{Where we stopped last time}
  \begin{columns}[T]
    \begin{column}{0.5\textwidth}
      A plain RNN rewrites its whole note every step:
      \[ \bm h_t = \tanh(W_{xh}\bm x_t + W_{hh}\bm h_{t-1} + \bm b) \]
      The gradient that reaches $k$ steps back has been multiplied by $W_{hh}$ and by $\tanh'$ $k$ times.
      \vskip0.6em
      In the \emph{Hidden State} game, the gradient survived 40 steps only when $|w_{hh}|$ was between $0.944$ and $1.059$.
    \end{column}
    \begin{column}{0.46\textwidth}
      \textbf{Who remembers this?}\\[0.3em]
      ``The \alert{keys} to the old wooden cabinet in my grandmother's house \underline{\hspace{1.2cm}} on the table.''
      \vskip0.6em
      To choose \emph{is} or \emph{are}, the network must carry ``keys is plural'' across nine words of distraction, then use it once, then drop it.
      \vskip0.6em
      Today: a cell that can \textbf{keep}, \textbf{write} and \textbf{show} on purpose.
    \end{column}
  \end{columns}
\end{frame}

\begin{frame}{Plan for today}
  \begin{tabular}{@{}r@{\hspace{1.2em}}p{9.4cm}@{\hspace{1em}}l@{}}
    \toprule
    \textbf{min} & \textbf{What we do} & \textbf{Worksheet} \\ \midrule
    0--5   & Recap: why plain RNNs forget & \\
    5--30  & The LSTM cell: three gates and a conveyor belt; one step by hand & Part A \\
    30--40 & GRU: the same idea with two gates & Part B \\
    40--50 & Bidirectional and stacked RNNs & \\
    50--62 & Case study: how far back can each cell remember? & Part C \\
    62--80 & Game: \emph{Gatekeeper} & Part D \\
    80--90 & Exit ticket; the Colab lab for home & Part E \\ \bottomrule
  \end{tabular}
  \vskip1em
  \margin{Reading: Goodfellow et al., \emph{Deep Learning}, Section 10.10. C.\ Olah, ``Understanding LSTM Networks'' (2015) is an excellent visual companion.}
\end{frame}

\sectionslide{Part 1}{The LSTM cell}

\begin{frame}{The idea before the equations}
  \begin{columns}[c]
    \begin{column}{0.5\textwidth}
      Give the network two things to carry:
      \begin{itemize}
        \item a \textbf{cell state} $\bm c_t$: long-term memory, a conveyor belt along the top of the cell that changes only in small, deliberate ways;
        \item a \textbf{hidden state} $\bm h_t$: what the cell chooses to show the outside world at this step.
      \end{itemize}
      \vskip0.5em
      Three \textbf{gates}, each a number between 0 and 1 per memory slot, decide:
      \begin{enumerate}
        \item \textbf{forget} $f$: how much of the old belt to keep;
        \item \textbf{input} $i$: how much of the new candidate to write;
        \item \textbf{output} $o$: how much of the belt to reveal.
      \end{enumerate}
    \end{column}
    \begin{column}{0.46\textwidth}
      \textbf{A notebook analogy}\\[0.3em]
      \renewcommand{\arraystretch}{1.25}
      \begin{tabular}{@{}ll@{}}
        cell state & the notebook \\
        forget gate & the eraser \\
        input gate & the pen \\
        output gate & which page you show \\
        hidden state & what you say aloud \\
      \end{tabular}
      \vskip0.6em
      \small A plain RNN has no eraser and no pen: it writes the whole page again from scratch at every word.
    \end{column}
  \end{columns}
\end{frame}

\analogyframe{The LSTM is a file in an office}
{A file moves from desk to desk. At each desk a clerk can remove old pages that no longer matter, add a new note, and decide which page to show at today's meeting. The file itself is never rewritten from scratch.}
{the file that travels between desks & \maps & the cell state $\bm c_t$ \pause\\
 removing outdated pages & \maps & forget gate $\bm f_t$ \pause\\
 adding a new note sheet & \maps & input gate $\bm i_t$ with candidate $\bm g_t$ \pause\\
 the page you show at the meeting & \maps & output gate $\bm o_t$, giving $\bm h_t$ \pause\\
 the old RNN: retyping the whole file at every desk & \maps & $\bm h_t=\tanh(W\ldots)$ every step \\}

\begin{frame}{Inside one LSTM cell}
  \centering
  \resizebox{!}{5.6cm}{%
  \begin{tikzpicture}[x=1.05cm,y=0.95cm]
    \draw[rulegrey,fill=paperwash,rounded corners=6pt] (0.3,-0.2) rectangle (8.2,4.1);
    % belts
    \draw[arr,line width=1.4pt,color=navy] (-0.6,3.4) node[left,font=\small,text=ink]{$\bm c_{t-1}$} -- (9.1,3.4) node[right,font=\small,text=ink]{$\bm c_t$};
    \draw[line width=0.8pt] (-0.6,0.3) node[left,font=\small]{$\bm h_{t-1}$} -- (5.9,0.3);
    \draw[line width=0.8pt] (0.9,-0.9) node[below,font=\small]{$\bm x_t$} -- (0.9,0.3);
    \fill (0.9,0.3) circle (1.5pt);
    % gate layers
    \node[gate,fill=white,minimum size=0.6cm,visible on=<2->] (f) at (1.7,1.5) {$\sigma$};
    \node[gate,fill=white,minimum size=0.6cm,visible on=<3->] (i) at (3.0,1.5) {$\sigma$};
    \node[gate,fill=white,minimum size=0.6cm,font=\tiny,visible on=<3->] (g) at (4.2,1.5) {tanh};
    \node[gate,fill=white,minimum size=0.6cm,visible on=<4->] (o) at (5.9,1.5) {$\sigma$};
    \foreach \n/\s in {f/2,i/3,g/3,o/4} { \draw[arr,visible on=<\s->] (\n |- 0,0.3) -- (\n); }
    \foreach \x/\s in {1.7/2,3.0/3,4.2/3} { \fill[visible on=<\s->] (\x,0.3) circle (1.5pt); }
    % ops
    \node[op,visible on=<2->] (xf) at (1.7,3.4) {$\times$};
    \node[op,visible on=<3->] (xi) at (4.2,2.5) {$\times$};
    \node[op,visible on=<3->] (pl) at (4.2,3.4) {$+$};
    \node[op,font=\tiny,minimum size=0.55cm,visible on=<4->] (th) at (7.0,2.5) {tanh};
    \node[op,visible on=<4->] (xo) at (7.0,1.5) {$\times$};
    \draw[arr,visible on=<2->] (f) -- (xf); \draw[arr,visible on=<3->] (g) -- (xi); \draw[arr,visible on=<3->] (i) |- (xi); \draw[arr,visible on=<3->] (xi) -- (pl);
    \draw[arr,visible on=<4->] (7.0,3.4) -- (th); \fill[visible on=<4->] (7.0,3.4) circle (1.5pt);
    \draw[arr,visible on=<4->] (th) -- (xo); \draw[arr,visible on=<4->] (o) -- (xo);
    \draw[arr,visible on=<4->] (xo) |- (9.1,0.3) node[right,font=\small]{$\bm h_t$};
    \draw[arr,visible on=<4->] (7.6,0.3) -- (7.6,4.6) node[above,font=\small]{$\bm h_t$};
    % labels
    \node[font=\scriptsize,text=brick,visible on=<2->] at (1.2,1.0) {forget};
    \node[font=\scriptsize,text=brick,visible on=<3->] at (2.5,1.0) {input};
    \node[font=\scriptsize,text=brick,visible on=<3->] at (4.2,0.95) {cand.};
    \node[font=\scriptsize,text=brick,visible on=<4->] at (6.45,1.0) {output};
    \node[font=\scriptsize,text=navy,anchor=south] at (3.0,3.5) {cell state: the conveyor belt};
  \end{tikzpicture}}
  \vskip0.2em
  \uncover<5->{\small Each gate looks at the same two things, $\bm h_{t-1}$ and $\bm x_t$, through its own weights. Only $\times$ and $+$ touch the belt.}
\end{frame}

\begin{frame}{The equations (Hochreiter \& Schmidhuber, 1997; forget gate: Gers et al., 2000)}
  \vskip-0.3em
  \begin{columns}[T]
    \begin{column}{0.52\textwidth}
      \begin{align*}
        \bm f_t &= \sigma\big(W_f[\bm h_{t-1},\bm x_t] + \bm b_f\big) && \text{forget}\\
        \bm i_t &= \sigma\big(W_i[\bm h_{t-1},\bm x_t] + \bm b_i\big) && \text{input}\\
        \bm g_t &= \tanh\big(W_g[\bm h_{t-1},\bm x_t] + \bm b_g\big) && \text{candidate}\\
        \bm o_t &= \sigma\big(W_o[\bm h_{t-1},\bm x_t] + \bm b_o\big) && \text{output}\\[0.3em]
        \onslide<2->{\bm c_t} &\onslide<2->{= \bm f_t\odot \bm c_{t-1} \;+\; \bm i_t\odot \bm g_t && \text{\alert{update the belt}}}\\
        \onslide<3->{\bm h_t} &\onslide<3->{= \bm o_t\odot \tanh(\bm c_t) && \text{what is shown}}
      \end{align*}
    \end{column}
    \begin{column}{0.45\textwidth}
      \begin{itemize}
        \item $[\bm h_{t-1},\bm x_t]$ means the two vectors stacked into one.
        \item $\sigma$ squashes to $(0,1)$: a gate is a dimmer switch, not an on/off switch, so it stays differentiable.
        \item $\odot$ is element by element: each memory slot has its own gates.
        \item The candidate uses $\tanh$ so it can write positive or negative values.
      \end{itemize}
      \vskip0.4em
      \small Read the belt update aloud: \emph{keep this much of the old, add this much of the new}.
    \end{column}
  \end{columns}
\end{frame}

\analogyframe{A gate is a fan regulator, not a light switch}
{A light switch is on or off. A fan regulator can be anywhere from 0 to 5. Every gate in an LSTM is a regulator, and each memory slot has its own set, turned by the network at every word.}
{regulator at 0 & \maps & gate value $0$: nothing passes \pause\\
 regulator at 5 & \maps & gate value $1$: everything passes \pause\\
 regulator at 3 & \maps & gate value $0.6$: part passes \pause\\
 a smooth knob, not a click & \maps & the sigmoid $\sigma$: differentiable, so it can be learned \\}

\begin{frame}{The gates, word by word}
  \small
  One memory slot tracks ``is the subject plural?'' in\\[0.2em]
  \emph{The keys to the old cabinet \underline{\ are\ } on the table.}
  \vskip0.6em
  \renewcommand{\arraystretch}{1.3}
  \begin{tabular}{@{}l c c c l@{}}
    \toprule
    \textbf{word} & $f$ \textbf{(keep)} & $i$ \textbf{(write)} & $o$ \textbf{(show)} & \textbf{what happens on the belt} \\ \midrule
    The       & low  & low  & low  & clear out the previous sentence's subject \\
    \pause
    keys      & --   & \alert{high} & low  & write ``plural'' \\
    \pause
    to the old cabinet & \alert{high} & low & low & hold it; ``cabinet'' is singular but must not overwrite \\
    \pause
    \underline{\ \ ?\ \ } & high & low & \alert{high} & reveal ``plural'' so the output picks \emph{are} \\
    \pause
    on the table. & low & low & low & done with it; free the slot \\ \bottomrule
  \end{tabular}
  \vskip0.8em
  These values are illustrative. A trained LSTM learns them from data; nobody programs them. Researchers have found individual LSTM units that behave much like this, for example units tracking quotes and line lengths in text (Karpathy, Johnson \& Fei-Fei, 2015).
\end{frame}

\begin{frame}{Why the gradient survives}
  \begin{columns}[T]
    \begin{column}{0.5\textwidth}
      Along the belt,
      \[ \bm c_t = \bm f_t\odot \bm c_{t-1} + \bm i_t\odot \bm g_t \]
      so (holding the gates fixed)
      \[ \frac{\partial \bm c_t}{\partial \bm c_{t-1}} = \operatorname{diag}(\bm f_t). \]
      No weight matrix, no $\tanh'$. Going $k$ steps back multiplies by $f$ values only.
      \vskip0.4em
      If the network learns $f\approx 1$ for a slot, its gradient passes back almost untouched: $0.99^{100}=0.37$, while $0.9^{100}$ would be $0.00003$.
    \end{column}
    \begin{column}{0.46\textwidth}
      \textbf{Plain RNN}: memory path is \emph{multiplicative}, the same matrix every step.\\[0.3em]
      \textbf{LSTM}: memory path is \emph{additive}, with a per-slot, per-step gate the network controls.
      \vskip0.6em
      \textbf{A practical detail that matters.} Initialise the forget-gate bias to a positive value (1 is common; Jozefowicz et al., 2015) so that a new network starts out \emph{remembering}. You will see in the case study what happens without it.
    \end{column}
  \end{columns}
\end{frame}

\analogyframe{A photocopy of a photocopy, or the original file?}
{Photocopy a page, then copy the copy, fifty times: the text fades away. Now send the original file along instead, with each desk only stamping a note on it: page one is still crisp at the fiftieth desk.}
{copy of a copy of a copy & \maps & plain RNN: multiply by $W_{hh}$ and $\tanh'$ every step \pause\\
 the original file, passed on & \maps & LSTM: $\bm c_t = \bm f_t\odot\bm c_{t-1}+\ldots$ with $f\approx1$ \pause\\
 a new clerk told ``keep everything unless told to remove it'' & \maps & forget-gate bias set positive at the start \\}

\begin{frame}{One LSTM step by hand}
  \small
  One memory slot. $c_{t-1}=0.8$. The pre-activations (the numbers inside $\sigma$ and $\tanh$) have already been computed:
  \[ z_f = 2.0,\qquad z_i = 0.0,\qquad z_g = 1.0,\qquad z_o = 1.0 \]
  \renewcommand{\arraystretch}{1.3}
  \begin{tabular}{@{}l l l@{}}
    \toprule
    forget & $f=\sigma(2.0)$ & $=0.881$ \\
    input & $i=\sigma(0.0)$ & $=0.500$ \\
    candidate & $g=\tanh(1.0)$ & $=0.762$ \\
    output & $o=\sigma(1.0)$ & $=0.731$ \\ \midrule
    \uncover<2->{belt & $c_t = 0.881(0.8) + 0.500(0.762) = 0.705+0.381$ & $=\alert{1.085}$} \\
    \uncover<3->{shown & $h_t = 0.731\,\tanh(1.085) = 0.731(0.795)$ & $=\alert{0.581}$} \\ \bottomrule
  \end{tabular}
  \vskip0.6em
  \uncover<3->{Note that $c_t$ went above 1. The belt is not squashed by $\tanh$; only what is \emph{shown} is.}
\end{frame}

\begin{frame}{Worksheet, Part A}
  \yourturn{5 min}
  \vskip1em
  \begin{block}{In pairs}
    \begin{enumerate}[<1->]
      \item The gates are given. Compute the new belt $c_t$ and the output $h_t$.
      \item Match three wishes (ignore a word, wipe the memory, stay silent) to gate settings.
    \end{enumerate}
  \end{block}
  \vskip0.6em
\end{frame}

\begin{frame}{What the gates cost}
  Each of the four blocks ($f, i, g, o$) is as big as a whole plain RNN layer:
  \[
    \text{RNN: } n(n+d)+n \qquad \text{GRU: } 3\big[n(n+d)+n\big] \qquad \text{LSTM: } 4\big[n(n+d)+n\big]
  \]
  \renewcommand{\arraystretch}{1.25}
  \begin{tabular}{@{}lrrr@{}}
    \toprule
    & \textbf{Plain RNN} & \textbf{GRU} & \textbf{LSTM} \\ \midrule
    $d=3$, $n=4$ & 32 & 96 & 128 \\
    $d=100$, $n=128$ (a word-embedding model) & 29{,}312 & 87{,}936 & 117{,}248 \\ \bottomrule
  \end{tabular}
  \vskip0.8em
  \small Libraries can differ slightly. Keras's GRU by default keeps two bias vectors per gate (\texttt{reset\_after=True}), giving $3[n(n+d)+2n]=88{,}320$ for the second row. Check with \texttt{model.summary()}.
\end{frame}

\sectionslide{Part 2}{The GRU: two gates instead of three}

\begin{frame}{Gated recurrent unit (Cho et al., 2014)}
  \begin{columns}[T]
    \begin{column}{0.5\textwidth}
      \begin{align*}
        \bm z_t &= \sigma\big(W_z[\bm h_{t-1},\bm x_t]+\bm b_z\big) && \text{update}\\
        \bm r_t &= \sigma\big(W_r[\bm h_{t-1},\bm x_t]+\bm b_r\big) && \text{reset}\\
        \tilde{\bm h}_t &= \tanh\big(W_h[\bm r_t\odot\bm h_{t-1},\,\bm x_t]+\bm b_h\big) && \text{candidate}\\[0.3em]
        \bm h_t &= (1-\bm z_t)\odot\bm h_{t-1} + \bm z_t\odot\tilde{\bm h}_t
      \end{align*}
      \vskip0.4em
      \small Conventions differ: PyTorch writes the last line with $\bm z_t$ and $1-\bm z_t$ swapped. Same model, the gate just means the opposite.
    \end{column}
    \begin{column}{0.46\textwidth}
      \textbf{What changed from LSTM}
      \begin{itemize}
        \item No separate belt: $\bm h_t$ is both memory and output.
        \item Forget and input are \emph{tied}: one gate $z$ decides the mix. Writing new information automatically erases the same amount of old.
        \item The reset gate $r$ decides how much of the past the candidate may look at.
        \item No output gate: the cell cannot hide what it remembers.
      \end{itemize}
      \vskip0.4em
      The memory path is still additive, so gradients survive the same way.
    \end{column}
  \end{columns}
\end{frame}

\analogyframe{GRU versus LSTM: a mixer tap versus two taps}
{A bathroom with two taps lets you run hot and cold independently, even both fully open. A single-lever mixer gives one control: turning towards hot takes away cold.}
{two separate taps & \maps & LSTM: forget $f$ and input $i$ set independently \pause\\
 both taps fully open, the bucket fills & \maps & $f=1,\ i=1$: the LSTM can add up, count \pause\\
 the single mixer lever & \maps & GRU: one update gate $z$ mixes old and new \pause\\
 more hot always means less cold & \maps & writing new always erases the same share of old \\}

\begin{frame}{Choosing between them}
  \small
  \renewcommand{\arraystretch}{1.3}
  \begin{tabular}{@{}>{\raggedright}p{3.1cm}>{\raggedright}p{3.2cm}>{\raggedright}p{3.2cm}>{\raggedright\arraybackslash}p{3.6cm}@{}}
    \toprule
    & \textbf{Plain RNN} & \textbf{GRU} & \textbf{LSTM} \\ \midrule
    Gates & none & update, reset & forget, input, output \\
    State carried & $\bm h$ & $\bm h$ & $\bm h$ and $\bm c$ \\
    Weights (same $n$) & $1\times$ & $3\times$ & $4\times$ \\
    Long-range memory & poor & good & good \\
    Can count / accumulate & no & weakly (writing erases) & yes (add with $f=1$, $i=1$) \\
    Reach for it when & short sequences, teaching & smaller data or compute budget & the default; long or structured sequences \\ \bottomrule
  \end{tabular}
  \vskip0.8em
  Large comparisons (Chung et al., 2014; Jozefowicz et al., 2015; Greff et al., 2017) found no clear winner between GRU and LSTM across tasks. Try both; tune the size and learning rate first.
\end{frame}

\begin{frame}{Worksheet, Part B}
  \yourturn{3 min}
  \vskip1em
  \begin{block}{GRU step}
    Previous state $h_{t-1}=0.6$, update gate $z=0.73$, candidate $\tilde h=0.20$.
    \begin{enumerate}[<1->]
      \item Compute $h_t=(1-z)h_{t-1}+z\tilde h$.
      \item Which $z$ keeps the old state? Which replaces it?
    \end{enumerate}
    \vskip0.4em Think about, no need to write: why can a GRU not hold 5 and then add 3 to get 8?
  \end{block}
\end{frame}

\sectionslide{Part 3}{Reading both ways, and stacking}

\analogyframe{Why read backwards too? Malayalam puts the verb last}
{``Avan oru chaaya \ldots'' \ He a tea \ldots\ drank? bought? spilled? In Malayalam you often learn what happened only at the last word. A good examiner, likewise, reads the whole answer before deciding what the first line meant.}
{the listener who waits for the verb & \maps & the backward RNN, reading right to left \pause\\
 understanding a word using what came after it & \maps & $\bm h_t=[\overrightarrow{\bm h}_t;\overleftarrow{\bm h}_t]$ \pause\\
 a live interpreter cannot wait for the end & \maps & no bidirectional model for streaming speech \\}
\begin{frame}{Bidirectional RNNs (Schuster \& Paliwal, 1997)}
  \begin{columns}[c]
    \begin{column}{0.52\textwidth}
      \centering
      \begin{tikzpicture}[x=1.3cm]
        \foreach \w [count=\i] in {Kerala,Blasters,won,today} {
          \node[font=\small\itshape] (w\i) at (\i,0) {\w};
          \node[cell,minimum width=0.85cm,minimum height=0.55cm] (f\i) at (\i,1.0) {};
          \node[cell,minimum width=0.85cm,minimum height=0.55cm,draw=brick,visible on=<2->] (b\i) at (\i,1.9) {};
          \node[io,minimum size=0.5cm,font=\tiny,visible on=<3->] (y\i) at (\i,3.0) {$y_\i$};
          \draw[arr] (w\i) -- (f\i); \draw[arr,visible on=<2->] (w\i.north east) to[bend right=30] (b\i.south east);
          \draw[arr,visible on=<3->] (b\i) -- (y\i); \draw[arr,visible on=<3->] (f\i.north west) to[bend left=30] (y\i.south west);
        }
        \foreach \i [evaluate=\i as \j using int(\i+1)] in {1,2,3} {
          \draw[arr,navy] (f\i) -- (f\j); \draw[arr,brick,visible on=<2->] (b\j) -- (b\i);
        }
        \node[font=\scriptsize,text=navy,anchor=east] at (0.5,1.0) {forward};
        \node[font=\scriptsize,text=brick,anchor=east,visible on=<2->] at (0.5,1.9) {backward};
      \end{tikzpicture}
    \end{column}
    \begin{column}{0.46\textwidth}
      Run one RNN left to right and a second, with its own weights, right to left. At each step join the two:
      \[ \bm h_t = [\,\overrightarrow{\bm h}_t\,;\,\overleftarrow{\bm h}_t\,] \]
      \textbf{Why.} Is ``Kerala'' a place or a team? In ``Kerala Blasters won'' you only know from the \emph{next} word. In ``Kerala received heavy rain'' it is a place.
      \vskip0.4em
      \textbf{Good for} tagging words, named entities, the encoder in translation, classifying a whole review.\\
      \textbf{Not for} live speech or next-word prediction: the future is not available yet.
    \end{column}
  \end{columns}
\end{frame}

\begin{frame}{Stacked (deep) recurrent layers}
  \begin{columns}[c]
    \begin{column}{0.45\textwidth}
      \centering
      \begin{tikzpicture}[x=1.15cm,y=0.95cm]
        \foreach \l in {1,2,3} { \foreach \t in {1,2,3,4} {
          \node[cell,minimum width=0.7cm,minimum height=0.5cm,visible on=<\l->] (n\l\t) at (\t,\l) {};
        } }
        \foreach \t in {1,2,3,4} {
          \node[io,minimum size=0.45cm,font=\tiny] (x\t) at (\t,0) {$x_\t$};
          \draw[arr] (x\t)--(n1\t); \draw[arr,visible on=<2->] (n1\t)--(n2\t); \draw[arr,visible on=<3->] (n2\t)--(n3\t);
        }
        \foreach \l in {1,2,3} { \foreach \t [evaluate=\t as \u using int(\t+1)] in {1,2,3} { \draw[rec,visible on=<\l->] (n\l\t)--(n\l\u); } }
        \foreach \l in {1,2,3} { \node[font=\scriptsize,anchor=east,text=inksoft,visible on=<\l->] at (0.45,\l) {layer \l}; }
      \end{tikzpicture}
    \end{column}
    \begin{column}{0.52\textwidth}
      The hidden states of layer 1 are the input sequence of layer 2.
      \begin{itemize}
        \item Lower layers tend to capture local patterns (characters, short phrases); higher layers longer structure.
        \item Two to four layers is typical. Beyond that, training gets hard without residual connections.
        \item Put dropout \emph{between} layers.
      \end{itemize}
      \vskip0.4em
      In Keras, every layer except the top one needs \texttt{return\_sequences=True}, so it hands the next layer a full sequence instead of only its last state.
    \end{column}
  \end{columns}
\end{frame}

\analogyframe{Stacked layers are an office hierarchy}
{The clerk reads every letter and notes the facts. The section officer reads the clerk's notes and spots patterns. The registrar reads the section officer's summary and sees the big picture.}
{the clerk, letter by letter & \maps & layer 1: characters, short patterns \pause\\
 the section officer, reading the clerk's notes & \maps & layer 2: phrases \pause\\
 the registrar & \maps & layer 3: the meaning of the whole sequence \\}

\sectionslide{Part 4 \textperiodcentered\ case study}{How far back can each cell remember?}

\begin{frame}{The experiment}
  \begin{columns}[T]
    \begin{column}{0.52\textwidth}
      \textbf{Task.} Each sequence starts with one of 8 ``key'' symbols, followed by $T-1$ random distractor symbols. At the end, name the key.
      \vskip0.3em
      {\centering\small\ttfamily \alert{K5} d3 d7 d1 d1 d6 \ldots\ d2 \quad $\rightarrow$ \quad K5\par}
      \vskip0.5em
      A coin-toss guess scores $1/8 = 12.5\%$. Nothing about the answer is in the last $T-1$ symbols, so the cell must hold one fact for $T$ steps.
      \vskip0.5em
      \textbf{Set-up.} 32 hidden units, Adam, learning rate $0.003$, gradient clipping at 5, 5{,}000 batches of 64, two random seeds. Test accuracy on 2{,}000 fresh sequences. Implemented in JAX; script in the course repository.
    \end{column}
    \begin{column}{0.45\textwidth}
      \textbf{Why this task?} It is the cleanest test of the claim from Part 1. In real data, a long dependency is mixed with many short ones and you cannot see which is which.
      \vskip0.5em
      \textbf{Gate initialisation.}\\ LSTM: forget-gate bias $+3$.\\ GRU: update-gate bias $-3$.\\ Both mean ``start out keeping the memory''.
      \vskip0.5em
      Before you see the results: \alert{write your prediction} in worksheet Part C. At what $T$ does each cell fail?
    \end{column}
  \end{columns}
\end{frame}

\begin{frame}{Results}
  \small
Test accuracy (\%) on 2{,}000 new sequences, one number per random seed. Chance is 12.5\%.
\vskip0.6em
\renewcommand{\arraystretch}{1.3}
\begin{tabular}{@{}l cccccc@{}}
  \toprule
  \textbf{sequence length} $T$ & 10 & 20 & 40 & 80 & 160 \\ \midrule
  Plain RNN & 100 / 100 & 100 / 100 & 100 / 100 & \alert{40 / 13} & \alert{27 / 12} \\
  \pause
  LSTM, forget bias $+3$ & -- & 100 / 100 & 100 / 100 & 100 / 100 & \alert{14} / 100 \\
  \pause
  GRU, update bias $-3$ & -- & 100 / 100 & 100 / 100 & 100 / 100 & 100 / 100 \\ \midrule
  \multicolumn{6}{@{}l}{\emph{First attempt, before we fixed the gate initialisation:}} \\
  \pause
  GRU, update bias 0 & 100 / 100 & 100 / 100 & \alert{13 / 11} & \alert{14 / 13} & \alert{12 / 12} \\ \bottomrule
\end{tabular}
\vskip0.8em
\begin{itemize}
  \item The plain RNN does better than its reputation up to $T=40$, then collapses to guessing.
  \item The gated cells hold the key for 80 steps, and the GRU for 160.
  \item One LSTM seed failed at $T=160$: training long memories is still not guaranteed.
\end{itemize}

\end{frame}

\begin{frame}{What we learned from it}
  \begin{columns}[T]
  \begin{column}{0.5\textwidth}
    \textbf{1. Gates move the limit, they do not remove it.} With the same training budget the plain RNN fails between $T=40$ and $T=80$; the gated cells reach $T=160$.
    \vskip0.6em
    \textbf{2. Initialisation can matter more than architecture.} With its update-gate bias at 0, the GRU starts by mixing half of the old memory away at every step: $0.5^{40}\approx 10^{-12}$. It never got a gradient to learn from and did \emph{worse} than the plain RNN.
  \end{column}
  \begin{column}{0.47\textwidth}
    \textbf{3. Report seeds, not single runs.} One LSTM seed at $T=160$ failed while the other was perfect. A single run would have told either story.
    \vskip0.6em
    \textbf{4. This is a toy.} On real text the gap is usually smaller, because most useful dependencies are short. Check it yourself with the IMDb lab.
    \vskip0.6em
    \margin{Script: \texttt{case\_study/memory\_length\_experiment.py}. About 10 minutes on a laptop CPU.}
  \end{column}
\end{columns}

\end{frame}

\begin{frame}[fragile]{Lab for home: sentiment on IMDb reviews}
  \small
  Run in Google Colab (free GPU: \emph{Runtime $\rightarrow$ Change runtime type}). About 3 minutes per model.
\begin{lstlisting}
from tensorflow import keras
from tensorflow.keras import layers

(x_tr, y_tr), (x_te, y_te) = keras.datasets.imdb.load_data(num_words=20000)
x_tr = keras.utils.pad_sequences(x_tr, maxlen=200)       # cut or pad every review to 200 words
x_te = keras.utils.pad_sequences(x_te, maxlen=200)

model = keras.Sequential([
    layers.Input(shape=(200,)),
    layers.Embedding(20000, 64),
    layers.Bidirectional(layers.LSTM(64)),                # swap: SimpleRNN, GRU, LSTM without Bidirectional
    layers.Dense(1, activation="sigmoid"),
])
model.compile(optimizer="adam", loss="binary_crossentropy", metrics=["accuracy"])
model.summary()                                           # compare the parameter counts with Part A
model.fit(x_tr, y_tr, batch_size=64, epochs=3, validation_split=0.2)
print(model.evaluate(x_te, y_te))
\end{lstlisting}
  Fill in the table in worksheet Part E for four versions of line 11. Bring it next class.
\end{frame}

\begin{frame}{Game: Gatekeeper}
  \yourturn{15 min}
  \begin{columns}[c]
    \begin{column}{0.36\textwidth}
      \centering
      \qrcode[height=4.0cm]{\gameurl gatekeeper.html}\\[0.4em]
      {\scriptsize\ttfamily \gameurlshort}
    \end{column}
    \begin{column}{0.6\textwidth}
      You are the gates of a one-slot LSTM. Tokens arrive one at a time. Before each one, open or close the \textbf{forget}, \textbf{input} and \textbf{output} gates. Then press \emph{Step} and watch the belt.
      \vskip0.5em
      Five levels: \emph{store and recall}, \emph{overwrite}, \emph{count}, \emph{recall twice}, and a last level where you play a \textbf{GRU} with only one gate.
      \vskip0.5em
      Wrong numbers on the output, or anything shown when it should be silent, costs points. Fill in worksheet Part D as you play.
    \end{column}
  \end{columns}
\end{frame}

\begin{frame}{What the game was teaching}
  \begin{itemize}
    \item \textbf{Forget means keep.} $f=1$ keeps the memory; $f=0$ erases it. Most people closed it on level 1 and lost the number.
    \item \textbf{Overwrite} needs two gates at once: $f=0$ to erase and $i=1$ to write.
    \item \textbf{Count} works because the belt adds: $f=1, i=1$ on every star.
    \item \textbf{Output} is separate from memory. The cell can remember without saying anything, which is what stops ``plural'' from leaking into every word.
    \item \textbf{GRU}: one gate $z$ moves memory and new input in opposite directions. Store and overwrite are easy; counting is impossible, because writing always erases.
  \end{itemize}
  \vskip0.6em
  In a trained network these gates are not 0 or 1 but values in between, learned by gradient descent, one per memory slot.
\end{frame}

\begin{frame}{Exit ticket (worksheet Part E)}
  \yourturn{5 min}
  \begin{enumerate}
    \item In an LSTM, setting the forget gate to 1 and the input gate to 0 makes the cell\\
      (a) erase its memory \quad (b) keep its memory unchanged \quad (c) copy the input \quad (d) output zero
    \item Why do LSTM gradients survive long gaps better than plain RNN gradients?\\
      (a) they use ReLU \quad (b) the cell state is updated additively, scaled only by the forget gate\\ (c) they have more parameters \quad (d) they are trained with a smaller learning rate
    \item Which task \emph{cannot} use a bidirectional LSTM?\\
      (a) named-entity tagging \quad (b) review sentiment \quad (c) live captioning of a speech as it is spoken \quad (d) POS tagging
  \end{enumerate}
  \vskip0.4em
  \uncover<+->{\small Answers: 1 (b) \quad 2 (b) \quad 3 (c)}
\end{frame}

\begin{frame}{To take away}
  \begin{enumerate}
    \item The LSTM keeps a cell state that changes only by gated multiplication and addition: $\bm c_t = \bm f_t\odot\bm c_{t-1}+\bm i_t\odot\bm g_t$.
    \item Forget, input and output gates are sigmoids in $(0,1)$, learned per memory slot and per step.
    \item The additive path lets gradients travel far back; initialise the forget bias so the cell starts by remembering.
    \item A GRU ties forget and input into one update gate, has no separate cell state, and uses 3 weight blocks instead of 4.
    \item Bidirectional RNNs read both ways when the whole sequence is available; stacked RNNs build hierarchy.
  \end{enumerate}
  \vskip0.8em
  \textbf{Coming up:} CNNs (Lectures 22--23), then reinforcement learning (Lecture 24), where an LSTM can supply the agent's state when one observation is not enough.
\end{frame}

\begin{frame}[plain]
  \begin{columns}[T]
    \begin{column}{0.63\textwidth}
      \vskip0.2em
      \begin{minipage}[c]{1.9cm}\includegraphics[width=1.9cm]{factsh_logo.png}\end{minipage}\hspace{0.8em}%
      \begin{minipage}[c]{0.66\linewidth}
        {\color{petrol}\factshfont\bfseries\LARGE FACTS-H Lab\par}
        \vskip0.2em
        {\scriptsize\color{inksoft}\textbf{F}air \& \textbf{A}ccountable \textbf{C}omputing Solutions that are \textbf{T}ransparent with \textbf{S}ociological and \textbf{H}umanistic Insights\par}
        \vskip0.2em
        {\scriptsize\color{inksoft}Indian Institute of Information Technology Kottayam\par}
      \end{minipage}
      \vskip0.4em
      {\small\itshape A research lab that builds tools people actually use, and publishes the science behind them.\par}
      \vskip0.3em
      \footnotesize
      \begin{itemize}[<1->]
        \item \textbf{Healthcare AI with AIIMS.} An ICMR-funded multicentric project on AI for ocular tumour detection; a clinical image-capture app used at four AIIMS sites; a gaze-tracking toolkit.
        \item \textbf{Software the institute runs on.} The attendance portal, Ph.D.\ lifecycle portal, exam seating automation, timetabling, grade-card verification, result analysis, and \emph{Aula}, a live polling tool for classrooms.
        \item \textbf{Three patents}, on eye tracking for film studies, nine-gaze capture of eye images, and secure authentication at the hostel gate.
        \item \textbf{People.} Twelve faculty and eleven Ph.D.\ scholars, among them Dr.\ Sara Renjit. The games you played today were built in the lab.
      \end{itemize}
    \end{column}
    \begin{column}{0.33\textwidth}
      \vskip1.2em
      \centering
      \qrcode[height=3.2cm]{\factshurl}\\[0.3em]
      {\scriptsize\ttfamily factsh.iiitkottayam.ac.in}
      \vskip1.2em
      \raggedright
      {\color{brass}\factshfont\bfseries\small Work with us\par}
      \vskip0.2em
      {\small You know your field's data and its hard problems. We know machine learning. If your department has a question that data could answer, talk to Dr.\ Sara Renjit after the session or write to \texttt{\scriptsize sararenjit@iiitkottayam.ac.in}.\par}
    \end{column}
  \end{columns}
  \begin{tikzpicture}[remember picture,overlay]
    \fill[petrol] ([yshift=0pt]current page.south west) rectangle ([yshift=4pt]current page.south east);
    \fill[brass] ([yshift=4pt]current page.south west) rectangle ([xshift=-0.75\paperwidth,yshift=6pt]current page.south east);
  \end{tikzpicture}
\end{frame}

\begin{frame}{References}
  \footnotesize
  \begin{itemize}[<1->]
    \item S.\ Hochreiter, J.\ Schmidhuber. Long short-term memory. \emph{Neural Computation} 9(8):1735--1780, 1997.
    \item F.\ A.\ Gers, J.\ Schmidhuber, F.\ Cummins. Learning to forget: continual prediction with LSTM. \emph{Neural Computation} 12(10):2451--2471, 2000.
    \item K.\ Cho et al. Learning phrase representations using RNN encoder--decoder for statistical machine translation. \emph{EMNLP} 2014.
    \item J.\ Chung, C.\ Gulcehre, K.\ Cho, Y.\ Bengio. Empirical evaluation of gated recurrent neural networks on sequence modeling. arXiv:1412.3555, 2014.
    \item R.\ Jozefowicz, W.\ Zaremba, I.\ Sutskever. An empirical exploration of recurrent network architectures. \emph{ICML} 2015.
    \item K.\ Greff et al. LSTM: a search space odyssey. \emph{IEEE Trans.\ Neural Networks and Learning Systems} 28(10), 2017.
    \item M.\ Schuster, K.\ K.\ Paliwal. Bidirectional recurrent neural networks. \emph{IEEE Trans.\ Signal Processing} 45(11), 1997.
    \item A.\ Karpathy, J.\ Johnson, L.\ Fei-Fei. Visualizing and understanding recurrent networks. arXiv:1506.02078, 2015.
    \item I.\ Goodfellow, Y.\ Bengio, A.\ Courville. \emph{Deep Learning}. MIT Press, 2016. Section 10.10.
  \end{itemize}
\end{frame}

\end{document}
